\section{Data and Empirical Design}

\subsection{Sample construction}

The daily panel spans January 2016 through June 2026 and begins from 88 futures products
across 13 broad economic sectors. Daily main contracts are selected causally using only
lagged activity. Returns are computed only when consecutive product observations refer to
the same underlying contract, preventing a change in quoted delivery contract from being
treated as an asset return. The analysis does not delete products or sectors because of
realized strategy performance.

\begin{table}[H]
\centering
\caption{Sanitized sample and design summary}
\label{tab:sample}
\begin{tabular}{lr}
\toprule
Item & Value \\
\midrule
Daily data span & 2016-01 to 2026-06 \\
Initial product set & 88 \\
Broad sectors & 13 \\
Evaluation period & 2021-01 to 2026-06 \\
Evaluation months & 66 \\
Mean eligible products per month & 65.9 \\
Momentum formation / skip & 11 months / 1 month \\
Reversal formation & 1 month \\
Long and short tails & 10\% each \\
Volatility history & 60 months \\
Router boundary & Fixed Q4 \\
\bottomrule
\end{tabular}
\end{table}

The distinction between the full data span and the strategy evaluation span is
important. Five years of observations are required before a monthly volatility can be
placed in its trailing distribution. The approximately ten-year panel therefore yields
only 66 evaluable holding months and 14 primary-proxy Q4 months. The effective tail-state
sample is small even though the underlying daily panel is broad.

\subsection{Sleeves}

For holding month $m$, the momentum score compounds product returns from $m-12$ through
$m-2$:
\begin{equation}
  MOM_{i,m}=\prod_{k=2}^{12}(1+R_{i,m-k})-1.
\end{equation}
The momentum sleeve buys the top decile and sells the bottom decile. The reversal score
is the most recent complete-month return,
\begin{equation}
  REVScore_{i,m}=R_{i,m-1},
\end{equation}
and the portfolio buys prior losers and sells prior winners. Each long and short leg is
equally weighted and has unit gross exposure.

\subsection{Volatility states and causal routing}

For each proxy, monthly market volatility is the sample standard deviation of daily
market returns within month $t$:
\begin{equation}
  \sigma_t=sd\{r^M_d:d\in t\}.
\end{equation}
The current observation is ranked in the trailing 60-month volatility distribution. The
state determined in $t$ first affects holdings in $t+1$:
\begin{equation}
R^{Router}_{t+1}=
\begin{cases}
R^{REV}_{t+1}, & \sigma_t\in Q4,\\
R^{MOM}_{t+1}, & \sigma_t\in\{Q1,Q2,Q3\}.
\end{cases}
\label{eq:router}
\end{equation}
The 75th-percentile boundary and both market proxies were specified without searching the
reported return sample.

\subsection{Evaluation logic}

We report arithmetic annual return, annualized volatility, Sharpe ratio, maximum
drawdown, conditional monthly means, and heteroskedasticity-and-autocorrelation-consistent
tests of return differences. Robust replication additionally requires that the Q4
relation be positive and representative across months, survive the equal-product proxy,
and improve on momentum and the static blend. Indicative implementation adjustments are
treated as a secondary feasibility check rather than a reason to modify signals.
